\scanpage{58}

Posons de plus
\[
(90)\qquad
\left\{\begin{aligned}
\nu'&=\nu_3-\nu_2-3\nu_2\nu_3+3\nu_0\\
\nu''&=\nu_3-\nu_2\unsure{\deleted{+}-|\deleted{-}+}2\nu_2\cdot\nu_3+2\nu_0.
\end{aligned}\right.
\]

On a alors les relations
\[
(91)\qquad
\left\{\begin{aligned}
\mu&=1-2\nu_2+4\nu'&&({}=\varepsilon_2+4\nu')\\
\mu&=1-2\nu_3+6\nu''&&({}=\varepsilon_3+6\nu'')\\[6pt]
&\cdot\\[-3pt]
\underbrace{-\nu_2+2\nu'}_{\nu}&=
\underbrace{-\nu_3+3\nu''}_{\nu}
&&\text{i.e. }2\nu'-\nu_2+\nu_3=3\nu''.
\end{aligned}\right.
\]
\margin{où $\varepsilon_p=(-1)^{\nu_p}$
\par d'où \deleted{$\unsure{\varepsilon_p=1-2\nu_p}$}
\par \uline{d'où} $\varepsilon_p-1=-2\nu_p$
\par $(p\in\{2,3\})$}

\uline{NB} La condition $\underline\mu=1$ équivaut aux conditions
\[
(92)\qquad \nu_2=\nu_3=\nu_0=0
\]
et implique
\[
(93)\qquad \nu=\nu'=\nu''=0
\]

\scanpage{59}
\mspage{677}

4. \unsure{Les groupes|Le groupe}
\deleted{$\widetilde Z_{\rho,\sigma}$ \unsure{et la suite} \illegible{2}}
$N_{\rho,\sigma}$, \deleted{\unsure{exacte}}\inserted{et la}
suite exacte
\diagram{\begin{tikzcd}[column sep=small]
1 \arrow[r] & N_{\rho,\sigma}^{0+} \arrow[r] & N_{\rho,\sigma} \arrow[r] & \unsure{\mathbb D}_{12}^{*} \arrow[r,dash] & 1.
\end{tikzcd}}
Relations avec $\unsure{\mathrm{Gl}|\mathrm{Sl}}(2,\mathbb Z/12\mathbb Z)$.

On le définit comme plongé dans
$\unsure{\mathbb D}_{12}^{*}\times\underline M\times\underline M$ :
\[
\begin{aligned}
(94)\quad&\widetilde Z_{\rho,\sigma}
\subset \unsure{\mathbb D}_{12}^{*}\times\underline M\times\underline M\\[-2pt]
&\prime\prime\\[-2pt]
&\bigl\{(\underline\mu,\alpha_0,\beta_0)\mid
(1-\rho_{\underline M})(\alpha_0)-(1-\sigma_{\underline M})(\beta_0)
=3\nu''(\underline\mu)\lambda_1\bigr\}
\end{aligned}
\]

Notons qu'on a
\[
(95)\qquad (1-\rho_{\underline M})
(\underbrace{\lambda_1-\lambda_0}_{\unsure{\delta}})=3\lambda_1
\]
donc la relation des cocycles dans (\unsure{77|72})
prend ainsi la forme
\[
(96)\qquad
(1-\rho_{\underline M})(\alpha_0-\nu''(\lambda_1-\lambda_0))
-(1-\sigma_{\underline M})(\beta_0)=0
\]
Je vais admettre \unsure{encore} qu'on a une structure de groupe naturelle sur
$\widetilde Z_{\rho,\sigma}$, pour laquelle la projection
$\widetilde Z_{\rho,\sigma}\to\unsure{\mathbb D}_{12}^{*}$
est morphisme de groupes. On trouve maintenant une suite exacte

(97)
\diagram{\begin{tikzcd}[column sep=large]
1 \arrow[r,dash] & \widetilde Z_{\rho,\sigma}^{+0\deleted{\unsure{*}}}
\arrow[r] & \widetilde Z_{\rho,\sigma} \arrow[r] & \unsure{\mathbb D}_{12}^{*} \arrow[r,dash] & 1
\end{tikzcd}}
où
\[
\begin{aligned}
(98)\quad&\widetilde Z_{\rho,\sigma}^{+0}
\overset{\mathrm{dfn}}=\operatorname{Ker}
(\widetilde Z_{\rho,\sigma}\mathbin{-}\unsure{\mathbb D}_{12}^{*})\\[-2pt]
&\rotatebox{90}{$\simeq$}\\[-2pt]
&\bigl\{(\alpha_0,\beta_0)\mid
(1-\rho_{\underline M})(\alpha_0)-(1-\sigma_{\underline M})(\beta_0)=0\bigr\}
\end{aligned}
\]
Ce schéma en groupes est lisse en car. $\ne3$, mais pas en car. $3$, où sa dim passe

\scanpage{60}

de la dim. relative $2$ à la dim relative $3$.
Ce fait d'une fibre exceptionnelle de dim. plus grande ne peut, il me semble, être
éliminé par une entourloupette, \deleted{\unsure{de}} en introduisant
\deleted{des} des définitions ad-hoc pour un ``meilleur'' groupe
$\widetilde Z_{\rho,\sigma}$ : il y a ``plus'' d'automorphismes de
$\underline{\unsure{\mathcal G}}_k$ (pour le $\unsure{\overline{\mathrm m}}$
corps alg. clos), ayant les propriétés envisagées, si $k$ est de car. $3$, que dans les autres cas de caractéristique.

Cependant il me semble qu'on peut construire un sous-schéma en groupes
\deleted{\unsure{de}} de $\widetilde Z_{\rho,\sigma}$, que j'ai envie de noter
$N_{\rho,\sigma}$, qui est lisse sur $S_0=\operatorname{Spec}(\mathbb Z)$ et
coïncide avec $\widetilde Z_{\rho,\sigma}$ au dessus de
\deleted{$\unsure{S_{0,3}}$}
$S_{0,\unsure{3|2}}\unsure{{}=}\operatorname{Spec}\mathbb Z[1/3]=S_0\setminus\{3\}$.
Pour le décrire, je note que dans l'équation des lacets\inserted{ (96)},
le \deleted{deuxième}\inserted{premier} est
\unsure{automatiquement|automatique} dans
\deleted{$\underline M$, \unsure{qui}}\inserted{le sous-module de $\underline M$ correspondant}
à la congruence
\[
c_0+c_1\equiv0\ (3)\qquad(\mathrm{cf}\ 62)
\]
et si on introduit le sous-\deleted{groupe}\inserted{$\mathbb Z$-module} d'indice

\scanpage{61}
\mspage{678}

$3$ correspondant $M_0$ de $M$
\[
\begin{aligned}
(99)\qquad&M_0\subset M\\[-2pt]
&\prime\prime\ \mathrm{dfn}\\[-2pt]
&\{c_0\lambda_0+c_1\lambda_1\mid c_0+c_1\equiv0\ (3)\},
\end{aligned}
\]
d'où un fibré vectoriel $\underline M_0$ sur $S_0$ ---
on a en fait un homomorphisme de schémas

(100)
\diagram{\begin{tikzcd}[column sep=large]
\widetilde Z_{\rho,\sigma} \arrow[r] & \underline M_0\\
(\underline\mu,\alpha_0,\beta_0) \arrow[r,mapsto] &
(1-\rho_{\underline M})
(\underbrace{\alpha_0+\nu''(\lambda_1-\lambda_0)}_{\widetilde\alpha_0})
+(1-\sigma_{\underline M})(\beta_0)
\end{tikzcd}}
qu'on peut expliciter ainsi : on a un hom de \uline{$\mathbb Z$-modules}

(101)
\diagram{\begin{tikzcd}[column sep=large]
M\times M\times\mathbb Z \arrow[r] & M_0\\
(\alpha_0,\beta_0,\nu'') \arrow[r,mapsto] &
(1-\rho_M)(\alpha_0+\nu''(\lambda_1-\lambda_0))
+(1-\sigma_M)(\unsure{\beta_0})
\end{tikzcd}}
d'où un hom \unsure{pour} les schémas en modules correspondants

(102)
\diagram{\begin{tikzcd}[column sep=large]
\underline M\times\underline M\times\underbrace{\underline{\mathbb Z}}_{\mathbb G_a}
\arrow[r] & \underline{M_0}
\end{tikzcd}}
qu'on compose avec

(103)
\diagram{\begin{tikzcd}[column sep=large]
\widetilde Z_{\rho,\sigma} \arrow[r] & \underline M\times\underline M\times\underline{\mathbb Z}\\
(\underline\mu,\alpha_0,\beta_0) \arrow[r,mapsto] & (\alpha_0,\beta_0,\nu''(\mu))
\end{tikzcd}}
On peut alors considérer l'équation des cocycles (96) comme équation \uline{dans}
$\underline M_0$ --- elle devient ainsi plus stricte,

\scanpage{62}

d'où un sous-schéma fermé de $\widetilde Z_{\rho,\sigma}$,
\[
\begin{aligned}
(104)\quad&N_{\rho,\sigma}\subset\widetilde Z_{\rho,\sigma}\\[-2pt]
&\prime\prime\\[-2pt]
&\{(\underline\mu,\alpha_0,\beta_0)\in\widetilde Z_{\rho,\sigma}\mid
\text{l'équation des \deleted{cocycles}\inserted{lacets}}\\
&\hspace{5cm}\text{(96) valable dans }\underline M_0\}.
\end{aligned}
\]
Considérons alors le morphisme
\diagram{\begin{tikzcd}
N_{\rho,\sigma} \arrow[r] & \unsure{\mathbb D}_{12}^{*}
\end{tikzcd}}
\margin{NB Le groupe $N_{\rho,\sigma}$ devrait pouvoir se noter aussi
$\unsure{\widetilde N}_{\rho,\sigma}[0]$, sauf erreur il
\par \unsure{Ou|En} plutôt ce serait
$(\unsure{M}_{\rho,\sigma}^{\unsure{b}}/\unsure{M}_{\rho})[0]$
\deleted{\unsure{d'après}} \unsure{avec} les notations du
$\S\,49$, \deleted{\unsure{2}}\inserted{\uline{$\unsure{\mathrm{VI}|\mathrm{IV}}$}}}
Il est clair que ce morphisme est encore un épi, la fibre d'un
$\underline\mu=(\nu_2,\nu_3,\nu_0)$ étant formée des
$(\alpha_0,\beta_0)\in\Gamma\,\underline M\times\underline M$
satisfaisant (96) dans $M_0$,
\deleted{On \unsure{prend le hom composé}}
\inserted{\unsure{s'identifie|s'identifiant} au \unsure{schéma} des}
$(\widetilde\alpha_0,\beta_0)$ satisfaisant l'équation linéaire
\diagram{\begin{tikzcd}
\deleted{M\times M\times\unsure{\mathbb Z}}
\arrow[r,phantom,"{\deleted{\longrightarrow}}" description]
& \deleted{M_0}
\end{tikzcd}}
en $\widetilde\alpha_0,\beta_0$
\[
(105)\qquad
(1-\rho_{\underline M})(\widetilde\alpha_0)
\unsure{-|=}(1-\sigma_{\underline M})(\beta_0)=0
\qquad\inserted{\text{dans }\underline M_0}.
\]
Or l'hom de Modules

(106)
\diagram{\begin{tikzcd}[column sep=large]
{} \arrow[d,phantom,"{\Bigg\{}" description] & \underline M\times\underline M \arrow[r] & \underline M_0\\
& (\widetilde\alpha_0,\beta_0) \arrow[r,mapsto] &
(1-\rho_{\underline M})(\widetilde\alpha_0)-(1-\sigma_{\underline M})(\beta_0)
\end{tikzcd}}
est un épi lisse --- étant associé à un hom
\diagram{\begin{tikzcd}
M\times M \arrow[r] & M_0
\end{tikzcd}}
correspondant qui est \unsure{manifestement|notoirement} épi.
Donc \deleted{$\unsure{\widetilde Z}$} $N_{\rho,\sigma}$ \uline{est lisse} sur $S_0$.
\deleted{\unsure{Ceci}} Étant un sous-schéma fermé de $\widetilde Z_{\rho,\sigma}$

\scanpage{63}
\mspage{\unsure{689|679}}

qui coïncide avec lui au-dessus de $S_0\setminus\{3\}$, il
s'ensuit que c'est l'\uline{adhérence schématique}
\unsure{que|par}
\[
\begin{gathered}
(107)\qquad
N_{\rho,\sigma}=\widetilde Z_{\rho,\sigma}|(\mathrm{Spec}(\mathbb Z)\setminus\{3\}),\\
\text{(adh. schématique)}
\end{gathered}
\]
ce qui implique que c'est aussi un sous-schéma en groupes.

C'est peut-être le moment de regarder d'un peu plus près le
sous-$\mathbb Z$-module $M_0$ de $M$,
\[
\begin{aligned}
(108)\qquad M_0
&=\mathbb Z(\lambda_1-\lambda_0)+\mathbb Z(3\lambda_1)\\
&=\mathbb Z\underbrace{(\lambda_1-\lambda_0)}_{\delta}
 +\mathbb Z\unsure{(1-\rho_M)}\underbrace{(\lambda_1-\lambda_0)}_{\delta}\\
&=\mathbb Z(\delta\unsure{,|)}\oplus\mathbb Z(1-\rho_M)(\delta)\\
&=\mathbb Z\delta+\mathbb Z\rho_M\delta.
\end{aligned}
\]
Notons que $M_0$ est stable sous $\rho$, \inserted{$\sigma$}, de façon
générale par construction, pour $\forall\ x\in M$, on
\[
\begin{aligned}
(109)\qquad x&\equiv\rho_M(x)\quad(M_0)
 &&\text{(i.e. }(1-\rho_M)(x)\in M_0\text{)}\\
x&\equiv\sigma_M(x)\quad(M_0)
 &&(1-\sigma_M)(x)\in M_0.
\end{aligned}
\]
\deleted{et}

Il est bien sur stable aussi par $\tau_M=-\mathrm{id}_M$.

Donc $M_0$ est \unsure{normalisé} par $\mathcal G$, donc
\inserted{en tant que sous-groupe de $\mathcal G$)} distingué dans
$\mathcal G$.
\deleted{C'est}\inserted{\uline{Il correspond donc à}} un sous-groupe
distingué de $\mathcal G_0=\mathrm{Gl}(2,\mathbb Z)$,

\scanpage{64}

contenant $[\Pi_0,\Pi_0]$, d'indice $3$ dans $\Pi_0$.
Je présume qu'il \deleted{est} \uline{contient} le sous-groupe de
congruences relatif à $12$ --- i.e. qu'il correspond à un sous-groupe
\inserted{distingué} de $\mathrm{Sl}(2,\mathbb Z/12\mathbb Z)$.
Il est d'indice $3.12=36$ dans $\mathrm{Sl}(2,\mathbb Z)$,
\deleted{$72$ dans $\mathrm{Gl}(2,\mathbb Z)$},
\deleted{\unsure{tandis}} que
\[
\mathrm{Sl}(2,\mathbb Z/12\mathbb Z)\simeq
\underbrace{\mathrm{Sl}(2,\mathbb Z/4\mathbb Z)}_{48}\times
\underbrace{\mathrm{Sl}(2,\mathbb Z/\unsure{3}\mathbb Z)}_{24}
\]
--- ce serait donc un sous-groupe \uline{\unsure{d'ordre}}
$32$ de $\mathrm{Sl}(2,\mathbb Z/12\mathbb Z)$,
\inserted{c'est évidemment le sous-groupe distingué}
engendré par les éléments \deleted{\illegible{1}} l'image de
\deleted{$[\sigma_\infty,\lambda_\infty]$},
\[
\left\{
\begin{array}{l}
(\sigma(\rho)\rho^{-1})^3=[\lambda_0,\lambda_1],\ \text{ou }[\sigma_\infty,\lambda_\infty]\\
\text{cf (7)}\\
\delta\overset{\mathrm{dfn}}{=}\lambda_1\lambda_0^{-1}
\ \deleted{\unsure{\in\Pi_0}}\\
\rho(\delta)
\end{array}
\right.
\]
ou encore le sous-groupe \inserted{distingué} engendré par le seul élément
\deleted{par les images dans $\mathrm{Sl}(2,\mathbb Z/12\mathbb Z)$ de)}
\[
\deleted{\delta=\lambda_1\lambda_0^{-1},\quad
\rho(\delta)=\lambda_\infty\lambda_1^{-1},\quad
\rho^2(\delta)=\lambda_0\lambda_\infty^{-1}}
\]
\[
(110)\qquad
\left\{
\begin{aligned}
\delta&=\lambda_1\lambda_0^{-1}=\ell_1\ell_0^{-1}
=\begin{pmatrix}1&0\\2&1\end{pmatrix}
 \begin{pmatrix}1&+2\\0&1\end{pmatrix}\\
&=\begin{pmatrix}1&+2\\2&5\end{pmatrix}
\end{aligned}
\right.
\]
La question est donc si ce sous-groupe distingué de
$\mathrm{Sl}(2,\mathbb Z/12\mathbb Z)$ est bel et bien d'ordre $32$,
donc d'indice $36$ --- son ordre pourrait être plus grand, son indice
plus petit. Déjà quel est l'ordre de $\delta$ ?

\scanpage{65}
\mspage{690}

mod $12$ ? On a
\[
\begin{array}{ll}
\deleted{\unsure{(111)}}&
\left\{\begin{aligned}
\delta^2&=\begin{pmatrix}1&2\\2&5\end{pmatrix}
\begin{pmatrix}1&2\\2&5\end{pmatrix}
=\begin{pmatrix}5&12\\12&29\end{pmatrix}
\equiv\begin{pmatrix}5&0\\0&5\end{pmatrix}\quad(12)\\
\delta^4&\equiv1\quad(12),
\end{aligned}\right.\\[-1.5em]
\text{d'où}&
\end{array}
\]
c'est déjà un bon point ! On a \deleted{d'ailleurs}
\inserted{$\overline{\mathrm m}$}
\[
(112)\qquad
\diagram{\left\{\begin{array}{lcl}
\delta\equiv1\quad(2)&&\delta\not\equiv1\ (4),\ \delta\not\equiv1\ (3)\\
&\Downarrow&\\
\delta^2\equiv1\quad(4)&&\delta^2\equiv-1\quad(3)
\end{array}\right.}
\]


Considérons
\[
(113)\qquad
\diagram{\begin{tikzcd}[column sep=small]
\mathrm{Sl}(2,\mathbb Z/12\mathbb Z)\simeq
\underbrace{\mathrm{Sl}(2,\mathbb Z/4\mathbb Z)}_{48}\times
\underbrace{\mathrm{Sl}(2,\mathbb Z/3\mathbb Z)}_{24}
\arrow[r,"u\times v"]&
\underbrace{\mathcal G_0^+/\Pi_0}_{12}\times
\underbrace{\mathcal G_3^+}_{3}
\end{tikzcd}}
\]
où \unsure{:}
\[
\diagram{\begin{tikzcd}
u:\mathrm{Sl}(2,\mathbb Z/4\mathbb Z)\arrow[r]&\mathcal G_0^+/\Pi_0
\end{tikzcd}}
\]
est l'hom canonique surjectif, dont le noyau est isomorphe à
$\deleted{\unsure{\mathbb Z}}\mu\times\mu$ --- les trois éléments
non nuls du noyau sont \unsure{juste} $\lambda_0,\lambda_1,\lambda_\infty$
permutés entre eux par $\rho$, et
$\dot\lambda_1\dot\lambda_0^{-1}=\dot\lambda_1\dot\lambda_0
=\dot\lambda_\infty$ et $\ldots$ i.e. \deleted{\unsure{les}}
les $\dot\lambda_i$ ne sont autres que $\dot\delta$ et ses permutés
par $\rho$ --- et
\[
\diagram{\begin{tikzcd}
v:\mathrm{Sl}(2,\mathbb Z/3\mathbb Z)\arrow[r]&\mathcal G_3
\end{tikzcd}}
\]
est l'hom composé
\[
\diagram{\begin{tikzcd}[row sep=tiny]
\mathrm{Sl}(2,\mathbb Z/3\mathbb Z)\arrow[r,"\text{épi}"]&
\unsure{A_4|G_4}\arrow[r,"\text{épi}"]
\arrow[d,phantom,"\cap" description]&\mathcal G_3^+\\
&\mathcal G_{\mathbb P^1(\mathbb Z/3\mathbb Z)}&
\end{tikzcd}}
\]
dont le noyau est d'ordre $8$. Le noyau de $u\times v$ est d'ordre
$4.8=32$, et il contient évidemment $\dot\delta$. Cela montre que
c'est bien le sous-groupe \inserted{distingué} engendré par $\delta$
(dont l'ordre à priori était \deleted{bel et} \uline{multiple} de $32$), et

\scanpage{66}

que \deleted{$\unsure{M_0}$}
\inserted{le sous-groupe envisagé de $\Pi_0\subset\mathrm{Sl}(2,\mathbb Z)$,
d'indice $3$ dans $\Pi_0$,}
est l'image inverse de ce sous-groupe \deleted{dans $M$} précédent de
$\mathrm{Sl}(2,\mathbb Z/12\mathbb Z)$.

Résumons les résultats obtenus :

\uline{Proposition} Considérons dans $\mathcal G_0^+=\mathrm{Sl}(2,\mathbb Z)$
la tour de sous-groupes invariants ($\overline{\mathrm m}$ dans $\mathcal G_0$)
\[
(114)\qquad
\diagram{\begin{tikzcd}[column sep=small,row sep=small]
\mathcal G_0^+&
\Pi=\Pi[2]\arrow[l,hook',"6"']\arrow[d,equal,"\mathrm{dfn}"]&
\Pi_0\arrow[l,hook',"2"']\arrow[d,equal]&
\Pi_{00}\ (\supset[\Pi_0,\Pi_0])\arrow[l,hook',"3"']\arrow[d,equal]\\
&\begin{array}{c}
\{u\in\mathcal G_0^+\mid u\equiv1\ (2)\}\\
\deleted{\Vert}\\
\deleted{\unsure{\mathfrak{sl}(2,\mathbb F_2)}\simeq
\{u\in\mathrm M_2(\mathbb F_2)\mid\mathrm{Tr}\,u=0\}}\\
\unsure{(i)|(1)}
\end{array}
&\begin{array}{l}
\text{Sous-groupe}\\
\text{engendré}\\
\text{par les }-\ell_i=\lambda_i
\end{array}
&\begin{array}{l}
\text{Sous-groupe }\uline{\text{inv.}}\\
\text{engendré par}\\
\deleted{\text{par les }\lambda_i,}\\
\deleted{[\lambda_0,\lambda_1]\ \text{et}}\\
\lambda_1\lambda_0^{-1}=\ell_1\ell_0^{-1}
\end{array}
\end{tikzcd}}
\]

Considérons d'autre part le groupe quotient
\[
(115)\qquad
\diagram{\begin{tikzcd}[row sep=small,column sep=small]
\mathrm{Sl}(2,\mathbb Z/12\mathbb Z)
\arrow[r,phantom,"\simeq" description]
\arrow[d,phantom,"\simeq" description]&
\mathcal G_0^+/\underbrace{\Pi[12]}_{\simeq\{u\in\mathcal G_0^+\mid u\equiv1\ (12)\}}\\
\mathrm{Sl}(2,\mathbb Z/4\mathbb Z)\times
\mathrm{Sl}(\unsure{2},\mathbb Z/\unsure{3}\mathbb Z),&
\end{tikzcd}}
\]
et la tour de sous-groupes suivants
\[
(116)\qquad
\diagram{\left\{\begin{array}{rll}
&H=\mathrm{Sl}(2,\mathbb Z/4\mathbb Z)\times\mathrm{Sl}(2,\mathbb Z/3\mathbb Z)
&\text{ordre }48\times24\\
6\ \cup&\hfill\cup\hfill&\\
&\Pi^\bullet=\mathfrak{sl}(2,\mathbb F_2)\times\mathrm{Sl}(2,\mathbb Z/3\mathbb Z)
&\text{ordre }8\times24\\
2\ \cup&\hfill\cup\hfill&\\
&\Pi_0^\bullet=\mathfrak{sl}(2,\mathbb F_2)'\times\mathrm{Sl}(2,\mathbb Z/3\mathbb Z)
&\text{ordre }4\times24\\
3\ \cup&\hfill\cup\hfill&\\
&\Pi_{00}^\bullet=\mathfrak{sl}(2,\mathbb F_2)'\times
\underbrace{\mathrm{Syl}_2}_{\simeq\deleted{\unsure{D_4}}D_4}
&\text{ordre }4\times8,
\end{array}\right.}
\]

\margin{\uline{NB} $(\mathcal G_0^+)_{\unsure{\mathrm{ab}}}\simeq
\mathbb Z/12\mathbb Z$ s'identifie à un
\deleted{$\unsure{(115)|(116)}\ \unsure{\mathrm{Sl}(2)}$}
quotient de $H/\Pi_{00}^\bullet$ par un ss-groupe (116) d'ordre $3$
--- soit encore le ss-groupe engendré par l'élément
$(-\dot\rho,1)$ de
\[
\underbrace{\mathrm{Sl}(2,\mathbb Z/4\mathbb Z)/\mathfrak{sl}(2,\mathbb F_2)'}_{\mathcal G_0^+/\Pi_0}
\times\mathrm{Sl}(2,\mathbb Z/3\mathbb Z)/\mathrm{Syl}_2.
\]}

Où \deleted{\unsure{on a}}
\[
(117)\qquad
\diagram{\left\{\begin{array}{ll}
\text{a)}&\mathfrak{sl}(2,\mathbb F_2)\simeq
\mathrm{Ker}\bigl(\mathrm{Sl}(2,\mathbb Z/4\mathbb Z)\longrightarrow
\mathrm{Sl}(2,\mathbb Z/2\mathbb Z)\bigr)\\[.5em]
\text{b)}&\mathfrak{sl}(2,\mathbb F_2)'\unsure{\subset}\mathfrak{sl}(2,\mathbb F_2),
\quad\begin{array}{l}
\text{sous-groupe distingué }\simeq\\
\deleted{\unsure{\mathbb F_3|\mathbb F_2}}\ \mathbb F_2\times\mathbb F_2\ ;\
\text{dont les trois él. }\ne0\\
\text{sont }\dot\lambda_0,\dot\lambda_1,\dot\lambda_\infty\
\text{ (images}\\
\text{dans }\mathrm{Sl}(2,\mathbb Z/4\mathbb Z)\text{) --- c'est}\\
\text{un supplémentaire de sa}\\
\text{sous-alg. de Lie }\mu\simeq\mathbb F_2\ \text{des matr.}\\
\text{\unsure{ctes|scal.}}
\end{array}\\[.5em]
\text{c)}&\underset{\substack{\Vert\\
\mathrm{Ker}(\mathrm{Sl}(2,\mathbb Z/3\mathbb Z)\to\mathcal G_3^+)}}{\mathrm{Syl}_2}
\subset\mathrm{Sl}(2,\mathbb Z/3\mathbb Z),\quad
\begin{array}{l}
\text{sous-groupe de Sylow, d'ordre }8\\
\text{(qui est distingué)}
\end{array}
\end{array}\right.}
\]

\scanpage{67}
\mspage{691}

Ceci posé, la tour de sous-groupes (114) est image inverse de la tour (116).

\deleted{\unsure{R.|Rm}}
Par acquit de conscience, je vais reprendre l'\unsure{esquisse|esquise}
de démonstration précédente. On note que l'hom
$\mathcal G_0^+\longrightarrow H$ applique les groupes figurant dans
(114) dans ceux de (116) --- pour $\Pi$ et $\Pi_0$, on sait
d'ailleurs depuis un moment que ce sont les images inverses de
$\Pi^\bullet$, $\Pi_0^\bullet$ ! $\unsure{\ast|\leftarrow}$
(Pour ceci, il \unsure{suffisait|suffirait} d'examiner
$\mathrm{Sl}(2,\mathbb Z/4\mathbb Z)$, au lieu de
$\mathrm{Sl}(2,\mathbb Z/12\mathbb Z)$ \ldots).

Ainsi l'image inverse de $\Pi_{00}^\bullet$ dans $\mathcal G_0^+$
contient $\Pi_{00}$. Comme les deux sont d'indice $36$, ces deux
groupes sont égaux, cqfd. (Le $\overline{\mathrm m}$
\unsure{argument} pouvait s'appliquer au cas de $\Pi_0$, $\Pi_0^\bullet$ \ldots)

\uline{Corollaire} L'homomorphisme canonique
\[
\diagram{\begin{tikzcd}[column sep=small]
\mathcal G_0^*=\mathrm{Gl}(2,\mathbb Z)\arrow[r]&
\mathrm{Gl}(2,\mathbb Z/12\mathbb Z)/
\bigl(\Pi_{00}^\bullet\deleted{\unsure{\simeq}}=
\mathfrak{sl}(2,\mathbb F_2)'\times\mathrm{Syl}_2\bigr)
\end{tikzcd}}
\]
\unsure{se factorise} par
\[
\mathcal G=\deleted{\mathcal G_0}\mathrm{Gl}(2,\mathbb Z)/[\Pi_0,\Pi_0]
\]
étudié dans le présent \S\,5\unsure{1|0}, et l'image inverse de
$\Pi_{00}^\bullet$ n'est autre que le sous-groupe $M_0$ de (99) (108).

\margin{\uline{NB}
Mais $\mathcal G_0^{\unsure{+|*}}\longrightarrow
\unsure{\mathrm{Sl}|\mathrm{Gl}}(2,\mathbb Z/12\mathbb Z)$
ne se factorise \uline{pas} par $\mathcal G^{\unsure{+|*}}$,
car s'il en était ainsi $\Pi_0^{\unsure{\bullet|\prime}}$
serait un \unsure{ss-}quotient de $\Pi_0/[\Pi_0,\Pi_0]$ et
donc serait commutatif --- \unsure{or tous}
\unsure{on voit} \unsure{deux|donc} \unsure{qui} \unsure{sont}
l'un \unsure{de|et} l'autre.}

\scanpage{68}

Les considérations précédentes rendent assez \unsure{naturel|naturelle} la considération du revêtement d'ordre $3$ de $U_{0,\overline{\mathbb Q}}$ \inserted{\deleted{\unsure{ou de} $M_{1,\unsure{0}}$}} correspondant au sous-groupe $\Pi_{00}$ de $\Pi_0\simeq\Pi_1(U_{0,\overline{\mathbb Q}})$.

\margin{c'est $\overline{\mathrm m}$ un \unsure{revêtement} de $M_{1,1}$\dots}

On vérifie du reste que c'est une courbe \uline{de genre $1$} -- si donc on s'avise de travailler avec $[\Pi_{0,0},\Pi_{0,0}]$ au lieu de $[\Pi_0,\Pi_0]$, on trouve un groupe extension de $\Pi_0/[\Pi_0,\Pi_0]\simeq M$, par un groupe de rang $2$ qui est ``elliptique'', et sur lequel $M$ va opérer de façon hautement non triviale -- par la suite il nous faudrait examiner cette action, et la mettre en relation avec l'action qu'on a \unsure{trouvée}, subordonnée à l'hom.
\diagram{\begin{tikzcd}
\mathrm{Sl}(2,\mathbb Z)^{\wedge}\arrow[r] & \mathrm{Sl}(2,\widehat{\mathbb Z}).
\end{tikzcd}}

\margin{c'est une courbe elliptique d'invariant $0$ (car elle est munie d'un \unsure{automorphisme} d'ordre $3$)}

\scanpage{69}
\mspage{692}

5) Les groupes $\underline M_{\rho,\sigma}$, $\underline M_{\rho,\sigma}^{\sharp}$ (via $\underline N_{\rho,\sigma}$).

À tout $\ell\in\Gamma\unsure{\underline M_0}$ correspond un autom intérieur $\operatorname{int}(\ell)$, qui est l'identité sur $\underline M_0$. On aura, tautologiquement
\[
(118)\qquad
\left\{\begin{aligned}
\operatorname{int}(\ell)(\rho)&=\alpha_{\unsure{0|\cdot}}(\rho)\\
\operatorname{int}(\ell)(\sigma)&=\beta(\sigma)
\end{aligned}\right.
\qquad\alpha=\beta=\ell
\]

Mais \deleted{on aura} aura-t-on pour autant la relation des lacets (relative à $\underline\mu=1$ donc $\nu''=0$)
\[
(1-\rho_M)(\underbrace{\alpha}_{\ell})=(1-\sigma_M)(\underbrace{\beta}_{\ell}).
\]
\[
\underset{\unsure{(*)|\alpha_1}}{\text{i.e.}}\qquad
\rho_M(\ell)=\sigma_M(\ell)\qquad ?
\]

Posant encore
\[
(119)\qquad\ell=c_0\lambda_0+c_1\lambda_1
\]
\deleted{on aura}
la relation $(*)$ s'écrit
\[
c_0\lambda_1+c_1\underbrace{\lambda_\infty}_{-\lambda_0-\lambda_1}
=c_0\lambda_1+c_1\lambda_0
\]
\[
\text{i.e.}\qquad-\unsure{c_1}\lambda_0+(c_0-c_1)\lambda_1
=c_1\lambda_0+c_0\lambda_1
\]
i.e. $c_1=0$ i.e. $\ell=c_0\lambda_0$. Donc c'est sous \unsure{réserve|réserves} qu'on se restreigne à considérer des $\ell\in\Gamma\underline L_0^{\unsure{\underline M_0|\underline M}}\unsure{\Rightarrow}$; \unsure{correspondants|correspondant}

\diagram{\begin{tikzcd}[column sep=small]
{(120)\quad\deleted{\unsure{\underline L_0}}\quad\ell=c_0\lambda_0\in
\underline L_{0\deleted{\unsure{,\rho,\sigma}}}^{\underline M}
\overset{\mathrm{dfn}}{=}\operatorname{Im}(}
& \underline{\mathbb Z}\arrow[r]\arrow[d,phantom,"\simeq" description]
& \underline M)\\
& \mathbb G_a\arrow[d,phantom,"\unsure{\ni}" description] &\\
& c_0\arrow[r,mapsto] & c_0\lambda_0
\end{tikzcd}}

On a un \deleted{\unsure{donc}} diagramme
\diagram{\begin{tikzcd}[column sep=large]
{(121)\quad\mathbb G_a\simeq\underline L_0^{\unsure{\underline M_0|\underline M}}}
\arrow[r]\arrow[d]
& \underline N_{\rho,\sigma} & \deleted{\unsure{\underline L_0^{\underline M}}}\\
\underline{\mathcal G}^{+} & &
\end{tikzcd}}

\scanpage{70}

l'image de
\diagram{\begin{tikzcd}
{(\unsure{124|122})\quad\underline L_0^{\unsure{\underline M_0|\underline M}}}\arrow[r]
& \underline N_{\rho,\sigma}
\end{tikzcd}}
est invariante -- l'hom $(121)$ étant en fait compatible avec l'action de $N_{\rho,\sigma}$ -- via $\underline\mu$ sur $\underline L_0^{M_0}\simeq\mathbb G_a$ (i.e. induite par l'opération de $\underline N_{\rho,\sigma}$ \unsure{de} sur\unsure{!| l'} $\underline{\mathcal G}^{+}\supset\underline M\supset\underline L_0^{\underline M}$), et \deleted{l'}\inserted{les deux} opérations de \deleted{$\unsure{N_\rho|A_\rho}$} $\underline L_0^{\underline M}$ sur $\underline{\mathcal G}^{+}$ déduites de $(121)$ étant \unsure{identique|identiques}. Donc d'après la \unsure{théorie|théor.} générale on \unsure{trouve|trouvera} un schéma en groupes $\underline M_{\rho,\sigma}$ et un diagramme de suites exactes

\diagram{\begin{tikzcd}[column sep=large,row sep=large]
{(123)\quad1}\arrow[r]
&\underline L_0^{\underline M}\arrow[r]\arrow[d]
&\underline N_{\rho,\sigma}\arrow[r,Rightarrow]\arrow[d]
&\overset{\inserted{\overset{\mathrm{dfn}}{=}\underline N_{\rho,\sigma}/\underline L_0^{\underline M}}}{\underline\Gamma_{\rho,\sigma}}\arrow[r]\arrow[d,Rightarrow]
&1\\
1\arrow[r]
&\underline{\mathcal G}^{+\unsure{a|\alpha}}\arrow[r]
&\underline M_{\rho,\sigma}\arrow[r,Rightarrow]
&\underline\Gamma_{\rho,\sigma}\arrow[r]
&1
\end{tikzcd}}

\margin{avec les notations du \S\,49 $\overline{\underline{\mathrm{VI}}}$, le groupe noté ici $\underline M_{\rho,\sigma}$ serait sans doute noté $\underline M_{\rho,\sigma}/\unsure{H_\rho|H_{\underline\Pi}}$, et $\underline M_{\rho,\sigma}^{\sharp}$ serait $\underline M_{\rho,\sigma}^{\sharp}/\unsure{H_\rho}$ (loc. cit. $8^\circ$)}

Remplaçant $\underline{\mathcal G}^{+}$ par $\underline M$, on trouve de $\overline{\mathrm m}$ $\underline M_{\rho,\sigma}^{\sharp}$, et

\diagram{\begin{tikzcd}[column sep=large,row sep=large]
{(124)}\arrow[d,phantom,"\Bigg\{" description]
&1\arrow[r,dash]
&\underline L_0^{\underline M}\arrow[r,dash]\arrow[d]
&\underline N_{\rho,\sigma}\arrow[r,dash]\arrow[d]
&\underline\Gamma_{\rho,\sigma}\arrow[r,dash]\arrow[d,Rightarrow]
&1\\
&1\arrow[r]
&\deleted{\unsure{\underline{\mathcal G}^{+}}}\underline M\arrow[r]
&\underline M_{\rho,\sigma}^{\sharp}\arrow[r]
&\underline\Gamma_{\rho,\sigma}\arrow[r]
&1
\end{tikzcd}}
,

d'ailleurs on a un hom de $(124)$ dans $(123)$ (comparer \S\,\uline{49} $\overline{\underline{\mathrm{VI}}}$ $5^\circ$\dots), et
\[
(125)\qquad M_{\rho,\sigma}^{\sharp}\subset M_{\rho,\sigma}\qquad\dots
\]

\scanpage{71}
\mspage{693}

Il faut préciser la structure de $\Gamma_{\rho,\sigma}$, en notant que $L_0^{\underline M}\subset N_{\rho,\sigma}^{0+}$, donc
\diagram{\begin{tikzcd}
{(\unsure{126|125})\quad1}\arrow[r]
&\deleted{\unsure{N_{\rho,\sigma}}}\underline\Gamma_{\rho,\sigma}^{0+}\arrow[r]
&\underline\Gamma_{\rho,\sigma}\arrow[r]
&\underline{\mathbb D}_{12}^{*}\arrow[r]
&1
\end{tikzcd}}
avec
\diagram{\begin{tikzcd}
{(127)\quad1}\arrow[r]
&\underline L_0^{M}\arrow[r,thick]
&\underline N_{\rho,\sigma}^{0+}\arrow[r]
&\underline\Gamma_{\rho,\sigma}^{0+}\arrow[r]
&1\\
&&\deleted{\unsure{\text{c'est }C_f|\text{c'est }G_f}}&&
\end{tikzcd}}

Considérons \deleted{la}\inserted{l'}\deleted{suite exacte}\inserted{hom surj.} de $\mathbb Z$-modules

\diagram{\begin{tikzcd}[column sep=small,row sep=small]
{(128)\quad\unsure{N_{0,\sigma}^{0+}|N_{\rho,\sigma}^{0+}}=
\deleted{\unsure{M\times M}}\operatorname{Ker}(}
&M\times M\arrow[r]
&M_0)\deleted{\unsure{\longrightarrow0}}\arrow[d,phantom,"\cap" description]\\
&&M
\end{tikzcd}}
\diagram{\begin{tikzcd}[column sep=small]
\deleted{\unsure{N}}\arrow[r,mapsto]
&\deleted{\unsure{\alpha_0,\beta_0}}\;\deleted{\unsure{(\alpha_0,\beta_0)}}
\end{tikzcd}}
\diagram{\begin{tikzcd}[column sep=large]
{(\underset{\substack{\Vert\\\inserted{a_0\lambda_0+a_1\lambda_1}}}{\alpha_0},
\underset{\substack{\Vert\\\inserted{b_0\lambda_0+b_1\lambda_1}}}{\beta_0})}\arrow[r]
& (1-\rho_M)(\alpha_0)-(1-\sigma_M)(\beta_0)
\end{tikzcd}}
\[
\underbrace{\hphantom{(1-\rho_M)(\alpha_0)-(1-\sigma_M)(\beta_0)}}
\]
\inserted{noyau formé des $(\alpha_0,\beta_0)$ tels que
\[
\begin{aligned}
c_0&\overset{\mathrm{dfn}}{=}a_0+a_1+(b_1-b_0)=0\\
c_1&\overset{\mathrm{dfn}}{=}2a_1-a_0-(b_1-b_0)=0
\end{aligned}
\qquad\bigg|\ \text{cf }(61)
\]}

Soit $\underline N^{0+}$ le noyau, de sorte qu'on a
\[
(129)\qquad\underline N_{\rho,\sigma}^{0+}\simeq\underline N^{0+}.
\]

\uline{Considérons} l'hom
\diagram{\begin{tikzcd}
{(130)}\arrow[d,phantom,"\Bigg\{" description]
&\mathbb Z\arrow[r]&N^{0+}\\
&n\arrow[r,mapsto]&(n\lambda_0,n\lambda_0)
\end{tikzcd}}

l'image \deleted{est} s'interprète comme \unsure{formé|formée} des systèmes $(a_0,a_1,b_0,b_1)$ tels que \deleted{$\unsure{a_0=b_0}$}
\[
a_1=b_1=0,\qquad a_0=b_0,\qquad\unsure{''}
\]
\deleted{d'où donc \unsure{un} \illegible{2}}
\inserted{i.e. on a suite exacte}
\diagram{\begin{tikzcd}
\deleted{\unsure{(131)}}&
0\arrow[r]&\mathbb Z\arrow[r]&N^{0+}\arrow[r,thick]&\mathbb Z^3\\
&&&(a_0,a_1,b_0,b_1)\arrow[r,mapsto]&(a_1,b_1,a_0-b_0)
\end{tikzcd}}
donc \deleted{\unsure{ce}} $N^{0+}/\mathbb Z$ est \deleted{\unsure{plat|libre}} $\mathbb Z$-module libre \inserted{de rang $1$},

\scanpage{72}

donc \unsure{une} suite exacte de $\mathbb Z$-modules
\diagram{\begin{tikzcd}[column sep=small,row sep=small]
{(131)\quad1}\arrow[r,dash]
&\mathbb Z\arrow[r]
&N^{0+}\arrow[r]
&\Gamma^{0+}\arrow[r]\arrow[d,phantom,"{\Vert\ \mathrm{dfn}}" description]
&0\\
&n\arrow[r,mapsto]&(n\lambda_0,n\lambda_0)
&N^{0+}/\operatorname{Im}\mathbb Z&
\end{tikzcd}}
avec $\Gamma^{0+}$ $\mathbb Z$-module libre de rg $1$.

On en conclut
\[
(132)\qquad\Gamma_{\rho,\sigma}^{0+}\simeq\underline\Gamma^{0+}
\qquad(\simeq\mathbb G_a)
\]

\deleted{\unsure{Donc}} Ainsi, tous les schémas en groupes qu'on vient de construire sont lisses sur $\operatorname{Spec}\mathbb Z$.

Peut-on définir une section plus-ou-moins naturelle de l'extension
(\inserted{$N_{\rho,\sigma}$ de $\Gamma_{\rho,\sigma}$ par $L_0^{\underline M}$, donc de}
$M_{\rho,\sigma}^{\sharp}$ de
$\underbrace{\Gamma_{\rho,\sigma}}_{\inserted{\substack{\text{lisse de dim}\\\text{relative }2}}}$
par $\underline M$ ?

Compte tenu de l'équation des lacets qui décrit $N_{\rho,\sigma}$, il semble qu'on \unsure{trouve} une telle section \unsure{dès|des} qu'on choisit un scindage de
\diagram{\begin{tikzcd}[column sep=small,row sep=large]
&0\arrow[r]
&\overset{\inserted{\mathrm{rg}\,2}}{N^{0+}}\arrow[r]
&M\times M\arrow[r]
&\overset{\inserted{\mathrm{rg}\,2}}{M_0}\arrow[r]
&0&\text{et de}\\
{(133)}\arrow[u,phantom,"\Bigg\{" description]
&0\arrow[r]
&\underset{\inserted{\mathrm{rg}\,1}}{\mathbb Z}\arrow[r]
&N^{0+}\arrow[r]
&\underset{\inserted{\mathrm{rg}\,1}}{\Gamma^{0+}}\arrow[r]
&0&
\end{tikzcd}}
(catégorie des $\mathbb Z$-modules). Mais il faut vérifier que les sections \textquotedblleft ensemblistes\textquotedblright{} de schémas que \unsure{ça} permet de définir sont multiplicatives\dots

\scanpage{73}
\mspage{694}

\(6^\circ)\) \(\underline Z_{\rho,\sigma}\) et
\(\underline G_{\rho,\sigma}=\mathbb D_{12}^{*}\mathbin{.}\underline{\mathcal G}^{+},\quad
\underline G_{\rho,\sigma}^{\sharp}\simeq\mathbb D_{12}^{*}\mathbin{.}\underline{\mathcal G}^{\sharp}
=\mathbb D_{12}^{*}\mathbin{.}\underline M\).

Une fois introduits \(\mathbb D_{12}^{*}\) et \(M_0\subset M\), il me
semble que le groupe \(\underline Z_{\rho,\sigma}\) (en lieu
et place de \(\underline{\widetilde Z}_{\rho,\sigma}\), dans lequel on tient
compte des \(\alpha,\beta\) et non seulement
de \(\xi,\eta\)) \unsure{redevient} sympathique:
\[
(134)\quad
\begin{gathered}
\underline Z_{\rho,\sigma}\subset
\underline{\operatorname{Aut}}(\underline{\mathcal G}^{+})\times\mathbb D_{12}^{*}
\\[-2pt]
\text{\(\prime\prime\)}\\[-2pt]
\left\{(u,\underline\mu)\ \left|\
\begin{array}{ll}
\text{a)}&u(\rho)=\xi\rho\quad\text{(loc.) }\underline{\mathcal G}_{\tau}^{\sharp}\text{-conjugué à }\rho,\text{ par un}\\
&\alpha\in\underline{\mathcal G}_{\tau}^{\sharp}\text{ tel que }\unsure{\chi}(\alpha)=\varepsilon_3(\underline\mu)\\
\text{b)}&u(\sigma)=\eta\sigma\quad\text{(loc.) }\underline{\mathcal G}_{\tau}^{\sharp}\unsure{+|-}\text{conjugué à }\sigma,\text{ par un}\\
&\beta\in\underline{\mathcal G}_{\tau}^{\sharp}\text{ tel que }\chi(\beta)=\varepsilon_2(\underline\mu)\\
\text{c)}&\text{équation des lacets}\\
&\xi_0=\eta_0(\lambda_{\unsure{1}}^3)^{\nu''}\quad\text{soit}\quad\xi_0=\eta_0+3\nu''\lambda_1\\
&\hfill\text{dans }\underline M_0
\end{array}\right.\right.
\end{gathered}
\]
où \(\nu''=\nu''(\underline\mu)\in\Gamma\mathbb G_a\)
(cf. \unsure{p.\ v.} p.~686), et \unsure{ci-dess.\(^{t}\)})
\[
(135)\qquad
\left\{\begin{aligned}
\xi_0&=\xi-\nu_3\lambda_1\in\Gamma\underline M_0
&\nu_3&=\nu_3(\underline\mu)\\
\eta_0&=\eta\omega^{\nu_2}\in\Gamma\underline M_0
&\nu_2&=\nu_2(\underline\mu)
\end{aligned}\right.
\]
\deleted{On a que} \inserted{On aura} \deleted{doit supposer} de plus que
\[
(136)\qquad
\left\{\begin{aligned}
&\xi_0\quad\text{\unsure{localement|local.\(^{mt}\)} dans l'image de }(1-\rho_{\underline M})(\underline M)\\
&\eta_0\quad\text{\unsure{localement|local.\(^{mt}\)} dans l'image }(1-\sigma_{\underline M})(\underline M)
\end{aligned}\right.
\]
et on aura
\[
(137)\quad
\begin{aligned}
\underline Z_{\rho,\sigma}\simeq\underline\Sigma_{\rho,\sigma}^{\mathrm{eff}}
\overset{\mathrm{dfn}}{=}
\bigl\{(\underline\mu,\xi_0,\eta_0)&\in\Gamma\mathbb D_{12}^{*}\times\underline M_0\times M_0\ \big|\\
1^\circ)\quad\xi_0&=\eta_0+\deleted{\unsure{\xi}}\nu''(3\lambda_1)\quad\text{dans }\underline M_0\\
&\text{où }\nu''=\nu''(\underline\mu)\\
2^\circ)\quad&\text{on a }(136)\quad\deleted{\unsure{\bigr\}}}\ \bigr\}.
\end{aligned}
\]
Or

\((1-\rho_{\underline M}):\)
\diagram{\begin{tikzcd}[column sep=large]
\underline M\arrow[r]&\underline M_0
\end{tikzcd}}

\((1-\sigma_{\underline M}):\)
\diagram{\begin{tikzcd}[column sep=large]
\underline M\arrow[r]&\underline{\mathbb Z}=\mathbb G_a\arrow[r]&\underline M_0\\
(\unsure{b_0}\lambda_0+b_1\lambda_1)\arrow[r]&(b_1-\underset{\deleted{\unsure{4|\lambda}}}{b_0})\,\deleted{\unsure{\lambda}}\arrow[r,mapsto]&\lambda\underbrace{(\lambda_1-\lambda_0)}_{\delta}=\lambda\delta
\end{tikzcd}}

\scanpage{74}

donc (136) \deleted{\unsure{s'écrit}} est une condition sur
\(\eta_0\) seulement, savoir

\((136\mathrm{bis})\quad\eta_0\in\Gamma\operatorname{Im}\bigl(\)
\diagram{\begin{tikzcd}[column sep=large]
\mathbb G_a\arrow[r]&\underline M_0\\
\lambda\arrow[r,mapsto]&\lambda\delta
\end{tikzcd}}
\(\bigr)=\unsure{\Gamma}\operatorname{Ker}\bigl(\)
\diagram{\begin{tikzcd}[column sep=large]
M_0\arrow[r]&\mathbb G_a\\
c_0\lambda_0+c_1\lambda_1\arrow[r]&c_0+c_1
\end{tikzcd}}
\(\bigr)\)

\deleted{Ceci dit, on trouve donc \unsure{comme}}
\inserted{Notons qu'on a}
\[
\deleted{(138)\qquad\underline\Sigma_{\rho,\sigma}^{\mathrm{eff}}\simeq}
\]
\[
(138)\qquad
\left\{\begin{aligned}
M_0&=\mathbb Z\underbrace{\delta}_{\lambda_1-\lambda_0}
\oplus\mathbb Z\mathbin{.}\underbrace{(3\lambda_1)}_{(1-\rho)\delta}
\qquad\text{d'où}\\
\underline M_0&=\deleted{\unsure{\mathbb G_a}}\ \underline{\mathbb Z\delta}
\times\underline{\mathbb Z\mathbin{.}(3\lambda_1)}\ .
\end{aligned}\right.
\]
et \deleted{par suite} on trouve

\((139)\)
\diagram{\begin{tikzcd}[column sep=large]
\underline\Sigma_{\rho,\sigma}^{\mathrm{eff}}\arrow[r,"\simeq"]&
\mathbb D_{12}^{*}\times\deleted{\unsure{\underline{\mathbb Z}\delta}}\ \deleted{\unsure{\mathbb G_a}}\ \underset{\unsure{\in|\simeq}\mathbb G_a}{\underline{\mathbb Z\delta}}\\
(\underline\mu,\xi_0,\unsure{\eta_0|\beta_0})\arrow[r,mapsto]&(\underline\mu,\unsure{\eta_0|\beta_0})
\end{tikzcd}}
.

Car \(\unsure{\xi_0|\zeta_0}\) se récupère avec \(\underline\mu,\unsure{\eta_0|\beta_0}\)
\deleted{\unsure{d'après}} l'équation des lacets
\[
\deleted{\boxed{\begin{aligned}
\unsure{\alpha_0|\xi_0}&=(1-\rho_{\underline M})^{-1}
\bigl(\unsure{\eta_0|\beta_0}+\nu''\underbrace{(3\lambda_1)}_{(1-\rho_{\underline M})(\delta)}\bigr)\\
&\deleted{\unsure{=\nu''\delta+(1-\rho_{\underline M})^{-1}\eta_0}}\\
&=\nu''\delta+\unsure{\xi_0|\eta_0}
\end{aligned}}}
\]

\uline{NB} Si on se donne par contre \(\xi_0\), on en tire
par l'équation des lacets et (138) la valeur de \(\unsure{\eta_0|\beta_0}\),
et celle de \(\nu''\) -- donc si on connaît
en plus \(\nu_2\) et \(\nu_3\), on en déduit
la valeur de \(2\nu_0\) donc (sur une
base où 2 est section régulière de \(\mathscr O_S\)) \(\nu_0\),
i.e. finalement \(\underline\mu\). Donc sur une

\scanpage{75}
\mspage{695}

base \(S\) où 2 est régulier, on trouve

\((140)\)
\diagram{\begin{tikzcd}[column sep=large]
\underline\Sigma_{\rho,\sigma}^{\mathrm{eff}}(S)\arrow[r,hook]&
\bigl(\underline M_0\deleted{\unsure{(S)}}\times\{0,1\}\deleted{\unsure{(S)}}\times\{0,1\}\bigr)(S)\\
(\underline\mu,\xi_0,\eta_0)\arrow[r,mapsto]&(\xi_0,\varepsilon_2(\underline\mu),\varepsilon_3(\underline\mu))
\end{tikzcd}}

l'image étant formée des systèmes
\((\xi_0,\unsure{\varepsilon_2|\xi_2},\unsure{\varepsilon_3|\xi_3})\) tels que, posant
\[
\deleted{(\unsure{141|140})}\qquad
\xi_0=\unsure{\eta_0|\beta_0}\delta+\nu''(3\lambda_1)
\qquad\eta_0,\nu''\in\deleted{\unsure{\Gamma}}\mathbb G_a(S)
\]
grâce à (138), on ait
\[
(141)\qquad
\left\{\begin{aligned}
&\nu''+\nu_2-\nu_3[+2\nu_2\nu_3]\quad
\begin{gathered}
\inserted{\text{divisible par 2}}\\[-2pt]
\deleted{\text{de la forme }2\nu_0}\\[-2pt]
\deleted{\text{divisible par}}
\end{gathered}\\
&\underbrace{(-5)^{\nu_2}(+5)^{\nu_3}+\unsure{6}(\nu''+\nu_2-\nu_3+2\nu_2\nu_3)}_{1-2\nu_3-2\nu_2\nu_3+6\nu''}\quad\text{inv.}
\end{aligned}\right.
\]

En tout état de cause on voit que
\(\underline\Sigma_{\rho,\sigma}^{\mathrm{eff}}\) donc aussi \(\underline Z_{\rho,\sigma}\) est lisse de
dim relative 2

\((142)\)
\diagram{\begin{tikzcd}[column sep=large]
1\arrow[r,dash]&\unsure{\underline Z}_{\rho,\sigma}^{+0}\arrow[r]\arrow[d,phantom,"\rotatebox{90}{$\simeq$}"']&
\underline Z_{\rho,\sigma}\arrow[r]&\mathbb D_{12}^{\unsure{*|\circ}}\arrow[r,dash]&1\\
&\mathbb G_a\simeq\underline{\mathbb Z\delta}
\end{tikzcd}}

(l'isom
\diagram{\begin{tikzcd}[column sep=large,row sep=small]
&&\inserted{\underline\Sigma_{\rho,\sigma}^{\mathrm{eff}+0}}\\
\mathbb G_a\arrow[r,phantom,"\simeq"]&\underline Z_{\rho,\sigma}^{+0}\arrow[ur,"\inserted{\simeq}"]&
\deleted{\unsure{\mathbb G_a}}\arrow[l,phantom,"\deleted{\unsure{\xleftarrow{\simeq}}}"']
\end{tikzcd}}
étant donné par
\diagram{\begin{tikzcd}[column sep=large]
\lambda\arrow[r,mapsto]&(\underline\mu=1,\xi_0\deleted{\unsure{=\eta_0}}=\lambda\delta,\eta_0=\xi_0=\lambda\delta)
\end{tikzcd}}
)

On \deleted{aura} posera
\[
(143)\qquad
\unsure{\underline L_0^{\sharp}|L_0^{\sharp}}\overset{\mathrm{dfn}}{=}
\{\lambda\mathbin{.}\lambda_0\mid\lambda\in\Gamma\mathbb G_a\}\subset\underline M
\]
et on définit

\((144)\)
\diagram{\begin{tikzcd}[column sep=large]
\underline L_0^{\sharp}\arrow[r]&\underline Z_{\rho,\sigma}^{\unsure{+0}}\arrow[r,"\simeq","\bigl("{description,pos=0,anchor=east}]&
\underline\Sigma_{\rho,\sigma}^{\mathrm{eff}\unsure{+0}}\bigr)\\
\underset{\prime\prime\lambda\mathbin{.}\lambda_0}{\ell}\arrow[r,mapsto]&
(\operatorname{int}(\ell),1)\arrow[r,mapsto]&
\begin{gathered}
(\deleted{\unsure{\mu}}1,(1-\rho_{\underline M})(\ell),(1-\sigma_{\underline M})(\ell))\\
=(1,-\lambda\delta,-\lambda\delta)
\end{gathered}
\end{tikzcd}}

\scanpage{76}

donc on a
\[
(145)\qquad\underline L_0^{\sharp}\simeq\underline Z_{\rho,\sigma}^{+0}
\]
on aura
\[
(146)\qquad\underline\gamma_{\rho,\sigma}\overset{\mathrm{dfn}}{=}
\underline Z_{\rho,\sigma}/\underline L_0^{\sharp}\simeq\underline{\mathbb D}_{12}^{*}\ .
\]
Si on construit maintenant \(G_{\rho,\sigma}\)

\deleted{\((\unsure{147})\)}

comme
\[
(147)\qquad\underline G_{\rho,\sigma}
=(\underline Z_{\rho,\sigma}\mathbin{.}\underline{\mathcal G}^{+})/\underline L_0^{\sharp}
\]
on trouve

\((148)\)
\diagram{\begin{tikzcd}[column sep=large,row sep=large]
1\arrow[r]&\underline L_0^{\sharp}\arrow[r]\arrow[d]&
\underline Z_{\rho,\sigma}\arrow[r]\arrow[d]&
\overbrace{\gamma_{\rho,\sigma}}^{\mathbb D_{12}^{*}}\arrow[r]\arrow[d,phantom,"\rotatebox{90}{$\simeq$}" description]&1\\
1\arrow[r,dash]&\underline{\mathcal G}^{+}\arrow[r]&G_{\rho,\sigma}\arrow[r]&\gamma_{\rho,\sigma}\arrow[r]&1
\end{tikzcd}}
.

D'ailleurs on a un scindage canonique
de \(\underline Z_{\rho,\sigma}\) sur \(\gamma_{\rho,\sigma}=\mathbb D_{12}^{*}\)

\((149)\)
\diagram{\begin{tikzcd}[column sep=large]
\mathbb D_{12}^{\unsure{*|\wedge}}\arrow[r]&\underline Z_{\rho,\sigma}\arrow[r,"\simeq"]&\underline\Sigma_{\rho,\sigma}^{\mathrm{eff}}\\
\underline\mu\arrow[r,mapsto]&\unsure{u_{\underline\mu}|\underline\mu}\arrow[r]&
(\underline\mu,\xi_0=\nu''(3\lambda_1),\eta_0=0)
\end{tikzcd}}

\margin{\uline{NB} On pose
\(G_{\rho,\sigma}^{\sharp}=\underbrace{\mathbb D_{12}^{*}}_{\gamma_{\rho,\sigma}}
\mathbin{.}\underset{\underline M}{\underline{\mathcal G}^{\sharp}}\subset G_{\rho,\sigma}\).}

d'où
\[
(150)\qquad
\left\{\begin{aligned}
\underline Z_{\rho,\sigma}&\simeq\mathbb D_{12}^{*}\mathbin{.}\underline L_0^{\sharp}\\
G_{\rho,\sigma}&\simeq\mathbb D_{12}^{*}\mathbin{.}\underline{\mathcal G}^{+}
\end{aligned}\right.
\]
L'opération de \(\mathbb D_{12}^{*}\) sur \(\underline L_0^{\sharp}\) se fait
via
\diagram{\begin{tikzcd}[column sep=large]
\mathbb D_{12}^{*}\arrow[r,dash]&\mathbb G_m
\end{tikzcd}}
et homothéties, celle
de \(\mathbb D_{12}^{*}\) sur \(\underline{\mathcal G}^{+}\) se fait par
\[
(151)\quad
\left\{\begin{aligned}
u_{\underline\mu}(\rho)&=\bigl[\deleted{\unsure{3\nu''|\nu_3}}\unsure{6\nu''|6\nu}\lambda_1\bigr]\rho
\quad\text{(}=\lambda_1^{\unsure{6\nu''|6\nu}}\rho\quad\text{en écritures \deleted{plus} exp.}\text{)}\\
u_{\underline\mu}(\sigma)&=\unsure{\omega^{\nu_2}}\sigma\\
\deleted{\unsure{u_{\underline\mu}}}\qquad u_{\underline\mu}(x)&=\mu x\qquad\text{si }x\in\Gamma\underline M\ .
\end{aligned}\right.
\]
